% =====================================================================
%  1bat-ud2-7.tex  —  HORA 7: Pressupostos integrats (producte de polinomis)
%  (sense preàmbul: s'inclou des de main.tex)
% =====================================================================
\horatitol{Sessió 7 — Pressupostos integrats}%
{Calcular el cost total d'un casal de joves (suma de polinomis) i la part
subvencionada (producte d'un nombre per un polinomi).}

\subsection{Idea clau}

A la sessió anterior vam sumar i restar pressupostos expressats com a polinomis.
Avui introduïm una nova operació: \textbf{multiplicar un polinomi per un nombre}.
Quan un ajuntament subvenciona la meitat del cost total, cada partida del
pressupost es divideix exactament per dos. Matemàticament, això és
$\num{0,5}\per T(x)$: cal distribuir el $\num{0,5}$ a \emph{cada terme} del
polinomi, un per un.

\begin{idea}
\textbf{Producte d'un nombre per un polinomi (propietat distributiva):}
\[
k\per(ax^2 + bx + c) = (k\per a)x^2 + (k\per b)x + (k\per c).
\]
El nombre $k$ multiplica \emph{tots i cadascun} dels termes del polinomi,
conservant la part literal ($x^2$, $x$, terme independent).

\medskip
\textbf{Compte amb el grau:} en multiplicar per un nombre, el grau del polinomi
\emph{no canvia}. Un polinomi de grau~2 continua sent de grau~2 un cop multiplicat
per $\num{0,5}$.
\end{idea}

\subsection{Exemples}

\begin{exemple}[subvenció del 50\% del cost d'un servei social]
Un servei de suport a famílies té un cost total de $T(x) = 40x + \num{1200}$,
on $x$ és el nombre de famílies ateses. L'administració subvenciona el
$50\%$ del cost. Calcula la part subvencionada $S(x) = \num{0,5}\per T(x)$.
\begin{align*}
S(x) &= \num{0,5}\per(40x + \num{1200}) \\
     &= (\num{0,5}\per 40)x + (\num{0,5}\per\num{1200}) \\
     &= 20x + 600.
\end{align*}
Interpretació: la subvenció cobreix $20\eur$ per família i $600\eur$ de costos
fixos.
\end{exemple}

\begin{exemple}[cost total i part subvencionada d'un casal de joves]
Un ajuntament vol obrir un casal de joves. El cost de personal és
$P(x) = \num{1500}x + \num{2000}$ i el cost de manteniment i materials és
$M(x) = 50x^2 + 300x$, on $x$ és el nombre de joves atesos.

\medskip
\textbf{Pas 1 — Cost total} (suma de polinomis, repàs de la sessió anterior):
\begin{align*}
T(x) &= P(x) + M(x) \\
     &= (\num{1500}x + \num{2000}) + (50x^2 + 300x) \\
     &= 50x^2 + (\num{1500}+300)x + \num{2000} \\
     &= 50x^2 + \num{1800}x + \num{2000}.
\end{align*}

\textbf{Pas 2 — Part subvencionada} (la meitat del cost total):
\begin{align*}
S(x) &= \num{0,5}\per T(x) \\
     &= \num{0,5}\per(50x^2 + \num{1800}x + \num{2000}) \\
     &= 25x^2 + 900x + \num{1000}.
\end{align*}
Observació: el terme $25x^2$ és la meitat de $50x^2$, el terme $900x$ és la
meitat de $\num{1800}x$ i el $\num{1000}$ és la meitat de $\num{2000}$. El
grau del polinomi (2) no ha canviat.
\end{exemple}

\subsection{Exercicis resolts}

\begin{exresolt}[part que paga l'ajuntament d'un programa de beques]
Un programa de beques menjador té un cost total de
$B(x) = 3x^2 + 80x + 500$ (en euros), on $x$ és el nombre de
beques concedides. L'ajuntament en paga el $60\%$.
Calcula el cost que assumeix l'ajuntament, $A(x) = \num{0,6}\per B(x)$.
\solucio
Distribuïm $\num{0,6}$ a cada terme del polinomi:
\begin{align*}
A(x) &= \num{0,6}\per(3x^2 + 80x + 500) \\
     &= (\num{0,6}\per 3)x^2 + (\num{0,6}\per 80)x + (\num{0,6}\per 500)
       &&\text{(distributiva)} \\
     &= \num{1,8}x^2 + 48x + 300.
\end{align*}
\resposta{$A(x) = \num{1,8}x^2 + 48x + 300\eur$}
\end{exresolt}

\begin{exresolt}[cost net d'un projecte comunitari]
Un projecte comunitari té un cost total de $T(x) = 50x^2 + \num{1800}x + \num{2000}$
i rep una subvenció equivalent a $S(x) = 25x^2 + 900x + \num{1000}$
(la meitat del cost, calculada a l'exemple anterior). Calcula el cost net
$N(x) = T(x) - S(x)$ que ha d'assumir el municipi.
\solucio
\begin{align*}
N(x) &= (50x^2 + \num{1800}x + \num{2000}) - (25x^2 + 900x + \num{1000}) \\
     &= 50x^2 + \num{1800}x + \num{2000} - 25x^2 - 900x - \num{1000}
       &&\text{(canvio tots els signes)} \\
     &= (50-25)x^2 + (\num{1800}-900)x + (\num{2000}-\num{1000}) \\
     &= 25x^2 + 900x + \num{1000}.
\end{align*}
Observació: el cost net coincideix amb la subvenció, cosa lògica: si es
paga exactament la meitat, l'altra meitat resta a càrrec del municipi.
\resposta{$N(x) = 25x^2 + 900x + \num{1000}\eur$}
\end{exresolt}

\subsection{Exercicis per fer}

\begin{exercicis}
\begin{exlist}
  \item Calcula els productes d'un nombre per un polinomi:
    \begin{enumerate}[label=\alph*), leftmargin=1.6em, topsep=2pt]
      \item $3\per(2x + 5)$
      \item $4\per(x^2 - 3x + 1)$
      \item $\num{0,5}\per(6x^2 + 10x - 8)$
    \end{enumerate}

  \item Una entitat gestiona un menjador social amb cost
        $M(x) = 15x + 800$ (en euros), on $x$ és el nombre
        d'àpats diaris. L'administració finança el $75\%$ del cost.
        Calcula la part finançada $F(x) = \num{0,75}\per M(x)$ i
        interpreta el coeficient de $x$ en el resultat.

  \item El casal de joves de l'Exemple~\theexemple{} rep finalment una
        subvenció del $40\%$ del cost total $T(x)=50x^2+\num{1800}x+\num{2000}$
        (en lloc del 50\%).
    \begin{enumerate}[label=\alph*), leftmargin=1.6em, topsep=2pt]
      \item Calcula la nova subvenció $S'(x) = \num{0,4}\per T(x)$.
      \item Calcula el cost net $N'(x) = T(x) - S'(x)$.
    \end{enumerate}

  \item Un ajuntament té el cost de personal del casal $P(x)=\num{1500}x+\num{2000}$
        i vol doblar el servei (atendre el doble de joves amb el mateix
        cost unitari). Expressa el nou cost de personal $P_2(x) = 2\per P(x)$
        i calcula quant costaria atendre $x=30$ joves.

  \item \textbf{(Compte amb el signe)} Calcula:
        \[
        -\num{0,5}\per(4x^2 - 6x + 10).
        \]
        Abans de resoldre-ho, indica quin signe tindrà cadascun dels tres
        termes del resultat. Comprova-ho un cop calculat.

  \item \textbf{(Interpretació)} Una fundació gestiona dos programes. El
        primer té un cost $A(x) = 2x^2 + 60x + 400$ i el segon,
        $B(x) = 2x^2 + 60x + 400$ (exactament igual).
        La fundació rep una subvenció única de $k\per(A(x)+B(x))$, on $k=\num{0,3}$.
        Calcula la subvenció total i explica, en paraules, per
        què el resultat és el mateix que calcular $\num{0,3}\per A(x)+\num{0,3}\per B(x)$.
\end{exlist}
\end{exercicis}

% ---------------------------------------------------------------------
\fullsolucions{Sessió 7}
\begin{solucions}
\begin{sollist}
  \item (a) $6x+15$ \quad (b) $4x^2-12x+4$ \quad (c) $3x^2+5x-4$
  \item $F(x)=\num{11,25}x+600$. El coeficient $\num{11,25}$ és el cost
        per àpat que assumeix l'administració (el $75\%$ de $15\eur$).
  \item (a) $S'(x)=20x^2+720x+800$ \quad
        (b) $N'(x)=30x^2+\num{1080}x+\num{1200}$
  \item $P_2(x)=\num{3000}x+\num{4000}$; \; per a $x=30$: $\num{94000}\eur$
  \item $-\num{0,5}\per(4x^2-6x+10) = -2x^2+3x-5$
        (tots els signes s'inverteixen en multiplicar per un nombre negatiu)
  \item $\num{0,3}\per(A(x)+B(x))=\num{0,3}\per(4x^2+120x+800)=\num{1,2}x^2+36x+240$.
        És el mateix que $\num{0,3}\per A(x)+\num{0,3}\per B(x)$ per la
        propietat distributiva: aplicar el mateix percentatge a la suma és
        equivalent a aplicar-lo a cada programa per separat i sumar-ho.
\end{sollist}
\end{solucions}
